\chapter{Optimisation}
\label{ch:optimisation}

\begin{tcolorbox}[feynbox={Sky}{BBg},title={Before and After}]
\textbf{Before this chapter you need:} quadratics and curvature
(\S\ref{sec:quadratics}), derivatives and the second-derivative test
(\S\ref{sec:derivative}--\S\ref{sec:secondderiv}), partial derivatives, the
gradient and the Hessian
(\S\ref{sec:partials}--\S\ref{sec:gradient}, \S\ref{sec:hessian}), norms
(\S\ref{sec:norms}), eigenvalues (\S\ref{sec:eigen}), expectation
(\S\ref{sec:expectation}), the law of large numbers (\S\ref{sec:lln}).\\
\textbf{After it you will understand:} what gradient descent actually does, why
mini-batching is not merely a memory convenience, how a learning rate interacts
with curvature, and why a minimax game does not converge in the way a
minimisation does --- the structural reason GAN training is hard.
\end{tcolorbox}

Every model in this thesis is produced by minimising something. This chapter
takes the gradient of Chapter~\ref{ch:calculus2} and builds the algorithms that
use it, ending with the two-player problem that the rest of the thesis lives
inside.

%% ─────────────────────────────────────────────────────────────────────────────
\section{Objective Functions}
\label{sec:objective}

\situation{
You are packing a suitcase. You want everything to fit, nothing to break, and the
weight under the airline limit. Before you can say whether one arrangement beats
another, you have to decide what ``better'' means --- and different answers give
different packings.
}

\problem{
``Train the model'' is not yet a computation. It becomes one only when there is a
single number to make small, and the choice of that number determines everything
the model will and will not learn.
}

\workedarith{
\textbf{Squared error.} Predictions $\hat{y} = (2, 5, 8)$ against truth
$y = (3, 5, 6)$:
\[
  L = (2-3)^2 + (5-5)^2 + (8-6)^2 = 1 + 0 + 4 = 5 .
\]

\textbf{Absolute error} on the same data:
\[
  L = |2-3| + |5-5| + |8-6| = 1 + 0 + 2 = 3 .
\]

\medskip
\textbf{They rank differently.} Compare two candidate predictors for
$y = (3,5,6)$:
\begin{itemize}
  \item Predictor A errs by $(0, 0, 3)$: squared $= 9$, absolute $= 3$.
  \item Predictor B errs by $(1, 1, 1)$: squared $= 3$, absolute $= 3$.
\end{itemize}
Squared error prefers B decisively; absolute error is indifferent. Squaring
punishes one large mistake more than several small ones, so the objective encodes
a judgement about which failures matter.

\medskip
\textbf{Why it matters here.} On heavy-tailed returns, squared error is dominated
by the rare enormous days --- the same leverage that broke the ARCH-LM test in
\S\ref{sec:archtest}. An objective is never neutral.
}

\formaldef{
An \textbf{objective function} (or \textbf{loss}, or \textbf{cost})
$L : \R^n \to \R$ maps parameters to a single number to be minimised. The problem
\[
  \theta^{*} = \arg\min_{\theta} L(\theta)
\]
asks for the parameters attaining the smallest value. Maximisation is the same
problem applied to $-L$.

$L$ typically averages over data: $L(\theta) = \frac{1}{N}\sum_{i=1}^{N}\ell(\theta; x_i)$,
an estimate of the expected loss $\E[\ell(\theta; X)]$.
}

\analogybreak{
The objective is a proxy for what you actually want, never the thing itself. This
thesis wants generators that reproduce stylised facts; what gets minimised is an
adversarial loss that mentions none of them. The gap between the two is why a
separate evaluation framework exists at all --- and why a model can improve its
training loss while getting worse at the job.
}

\whyhere{
The three models minimise different objectives --- TimeGAN a sum of
reconstruction, supervised and adversarial terms; QuantGAN and FinGAN a WGAN-GP
critic loss. None of them optimises Wasserstein distance to the real marginal, ACF
error, or tail index. The evaluation metrics are deliberately \emph{not} the
training objectives, which is what stops the comparison being circular.
}

%% ─────────────────────────────────────────────────────────────────────────────
\section{Local versus Global Minima}
\label{sec:minima}

\situation{
You are dropped somewhere in a mountain range in thick fog and told to find the
lowest point. You can feel the slope underfoot. Walking downhill until the ground
is level lands you at the bottom of \emph{a} valley --- with no way of knowing
whether a deeper one lies over the ridge.
}

\problem{
Every practical algorithm is a fog-walker. It sees only local information, so it
can certify only local optimality. Claims about global optima require either a
special structure or an exhaustive search, and neither is available here.
}

\workedarith{
Take $f(x) = x^4 - 4x^2 + x$. Then
\[
  f'(x) = 4x^3 - 8x + 1 .
\]
Stationary points occur near $x \approx -1.44$, $x \approx 0.13$ and
$x \approx 1.32$. Evaluate:

\medskip
\begin{tabular}{@{}lrr@{}}
\toprule
$x$ & $f(x)$ & Type \\
\midrule
$-1.44$ & $-5.68$ & local minimum (the global one) \\
$0.13$  & $+0.06$ & local maximum \\
$1.32$  & $-3.52$ & local minimum \\
\bottomrule
\end{tabular}

\medskip
\textbf{The consequence.} Start the walk at $x = 2$ and the slope carries you to
$1.32$, where $f = -3.52$. You stop, correctly, at a point where $f'(x) = 0$ and
$f''(x) > 0$. And you are $2.16$ above the true minimum, with nothing in the local
information to tell you so.

\medskip
\textbf{Starting point decides the outcome.} From $x = -2$ the same algorithm
reaches $-1.44$ and the global minimum. Same objective, same rule, different
answer --- determined entirely by initialisation.
}

\formaldef{
$\theta^{*}$ is a \textbf{local minimum} if $L(\theta^{*}) \leq L(\theta)$ for all
$\theta$ in some neighbourhood of $\theta^{*}$, and a \textbf{global minimum} if
the inequality holds for all $\theta$ in the domain.

First-order condition: $\nabla L(\theta^{*}) = \mathbf{0}$. Second-order
sufficient condition: the Hessian is positive definite --- every eigenvalue
positive (\S\ref{sec:eigen}) --- which is the many-variable version of
$f''(a) > 0$.

A point with $\nabla L = \mathbf{0}$ and mixed-sign eigenvalues is a
\textbf{saddle}.
}

\analogybreak{
In high dimensions the intuition from one variable misleads. For a random function
of $n$ variables, a stationary point is a local minimum only if all $n$
eigenvalues happen to be positive, which becomes vanishingly unlikely as $n$
grows. Saddles, not spurious minima, are the dominant obstacle --- and a saddle
has a downhill direction, so an algorithm that is not exactly stationary can
eventually escape. That is part of why training very large models works better
than the one-variable picture predicts.
}

\whyhere{
Every result in this thesis is a local optimum, and nothing in the procedure
certifies more. This is also why seeds are reported per seed and never averaged
away: different seeds mean different initialisations, hence different valleys, and
the spread across seeds is a genuine part of the finding rather than noise to be
smoothed.
}

%% ─────────────────────────────────────────────────────────────────────────────
\section{Convexity}
\label{sec:convexity}

\situation{
A smooth bowl has one lowest point, and a marble released anywhere inside reaches
it. An egg carton has many, and where the marble stops depends on where you
dropped it. Whether a problem is easy or hard is largely the question of which
shape you have.
}

\problem{
The previous section showed local search cannot certify a global answer. There is
one important exception, and knowing whether you are inside it tells you how much
to trust the result.
}

\workedarith{
\textbf{The chord test.} A function is convex if the straight line joining any two
points on its graph never dips below the graph.

For $f(x) = x^2$, take $x = 1$ and $x = 3$, with $f = 1$ and $9$. The midpoint of
the chord is at height
\[
  \tfrac12(1 + 9) = 5 ,
\]
while the function at the midpoint $x = 2$ is $f(2) = 4$. Since $4 \leq 5$, the
chord lies above. \checkmark

\medskip
\textbf{A non-convex case.} $g(x) = x^4 - 4x^2 + x$ from \S\ref{sec:minima}, at
$x = -1.44$ and $x = 1.32$: heights $-5.68$ and $-3.52$, chord midpoint
$\tfrac12(-5.68 - 3.52) = -4.60$. The midpoint is $x = -0.06$, where
\[
  g(-0.06) = 0.000013 - 0.0144 - 0.06 = -0.074 .
\]
Here $-0.074 > -4.60$: the function lies \emph{above} the chord. Not convex.
\checkmark

\medskip
\textbf{Via the second derivative.} $f''(x) = 2 > 0$ everywhere for $x^2$ ---
convex. For $g$, $g''(x) = 12x^2 - 8$, which is negative for
$|x| < \sqrt{2/3} = 0.816$ --- so $g$ is concave in that band and convex outside.
Mixed sign, not convex. \checkmark
}

\formaldef{
$f$ is \textbf{convex} on a convex set if for all $x, y$ and $t \in [0,1]$
\[
  f\bigl(tx + (1-t)y\bigr) \leq t f(x) + (1-t) f(y) .
\]
For twice-differentiable $f$ this is equivalent to the Hessian being positive
semi-definite everywhere.

The essential consequence: \textbf{for a convex function every local minimum is
global}, and the set of minimisers is convex.
}

\analogybreak{
Convexity is the exception in this field. Squared error in a \emph{linear} model
is convex; put the same loss behind a neural network and it is not, because the
composition of a convex function with a non-linear map need not be convex. Every
objective minimised in this thesis is non-convex, so no global guarantee is
available anywhere in it.
}

\whyhere{
The reconstruction loss in TimeGAN's Phase 1 is quadratic in the reconstruction
error, and its bowl shape is why an $R^2$ of $+0.995$ is meaningful --- the
autoencoder really did approach the bottom of a well-behaved objective. Phase 3 is
adversarial and has no such structure, which is exactly the difference
\S\ref{sec:minimax} describes.
}

%% ─────────────────────────────────────────────────────────────────────────────
\section{Gradient Descent}
\label{sec:gradientdescent}

\situation{
The fog-walker again, now with a method: feel which way is steepest downhill, take
a step that way, and repeat. It requires nothing but local slope, which is
precisely what is available.
}

\problem{
Section~\ref{sec:gradient} established that $-\nabla L$ is the steepest downhill
direction. Turning that into an algorithm requires deciding how far to step, and
knowing when the procedure can be trusted to make progress.
}

\workedarith{
Minimise $f(x) = x^2$ from $x_0 = 4$, with $f'(x) = 2x$ and step size
$\eta = 0.1$. The rule is $x \leftarrow x - \eta f'(x)$:

\medskip
\begin{tabular}{@{}lrrr@{}}
\toprule
Step & $x$ & $f'(x)$ & next $x$ \\
\midrule
0 & $4.000$ & $8.000$ & $3.200$ \\
1 & $3.200$ & $6.400$ & $2.560$ \\
2 & $2.560$ & $5.120$ & $2.048$ \\
3 & $2.048$ & $4.096$ & $1.638$ \\
\bottomrule
\end{tabular}

\medskip
Each step multiplies $x$ by $1 - 2\eta = 0.8$. After $n$ steps
$x_n = 4(0.8)^n$, converging to $0$ geometrically. At $n = 30$:
$4(0.8)^{30} = 0.005$.

\medskip
\textbf{Two variables.} Minimise $f(x,y) = x^2 + 3y^2$ from $(2, 2)$ with
$\eta = 0.1$. Here $\nabla f = [2x, 6y]$:
\[
  x: 2 \to 2 - 0.1(4) = 1.6 , \qquad y: 2 \to 2 - 0.1(12) = 0.8 .
\]
The $y$ coordinate moves much faster, because the surface is steeper in $y$ ---
the elongated bowl of \S\ref{sec:hessian}, with Hessian $\mathrm{diag}(2,6)$.

Next step: $x \to 1.28$, $y \to 0.32$. The $x$ direction shrinks by $0.8$ per
step, $y$ by $0.4$. One global step size, two very different effective rates.
}

\formaldef{
\textbf{Gradient descent} iterates
\[
  \theta_{t+1} = \theta_t - \eta\,\nabla L(\theta_t) ,
\]
with \textbf{learning rate} $\eta > 0$.

For an $L$-smooth convex objective --- gradient Lipschitz with constant $L$ ---
convergence is guaranteed when $\eta < 2/L$, at rate $O(1/t)$. For non-convex
objectives the guarantee weakens to convergence toward a stationary point, with no
claim about which one.
}

\analogybreak{
Gradient descent is greedy: it takes the locally steepest direction, which is
rarely the direction pointing at the minimum. In the elongated bowl above it
zig-zags, making rapid progress across the narrow direction and crawling along the
shallow one. The number governing how bad this gets is the ratio of largest to
smallest Hessian eigenvalue --- the \textbf{condition number} --- which here is
$6/2 = 3$, mild. Real loss surfaces reach condition numbers in the thousands.
}

\whyhere{
Every parameter in every model in this thesis is updated by a variant of this
rule. The gradient-norm clip of $\texttt{max\_norm}=1.0$ applied in all three
models rescales $\nabla L$ whenever $\norm{\nabla L}_2 > 1$, bounding the step
length regardless of how steep the surface becomes.
}

%% ─────────────────────────────────────────────────────────────────────────────
\section{The Learning Rate}
\label{sec:learningrate}

\situation{
Descending a staircase in the dark. Tiny shuffles get you down safely but take all
night. Enormous strides are fast until one overshoots the step and you land
further down than intended --- or fall.
}

\problem{
The step size is the single most consequential setting in the algorithm, and it
has no universally right value: too small wastes computation, too large diverges,
and the boundary between them depends on curvature that changes as you move.
}

\workedarith{
Minimise $f(x) = x^2$ from $x_0 = 4$. Each step multiplies $x$ by $(1 - 2\eta)$.

\medskip
\begin{tabular}{@{}lrl@{}}
\toprule
$\eta$ & factor $|1-2\eta|$ & behaviour \\
\midrule
$0.01$ & $0.98$ & converges, slowly \\
$0.1$  & $0.80$ & converges well \\
$0.5$  & $0.00$ & reaches the minimum in one step \\
$0.9$  & $0.80$ & converges, oscillating in sign \\
$1.0$  & $1.00$ & oscillates for ever between $4$ and $-4$ \\
$1.5$  & $2.00$ & \textbf{diverges}: $4 \to -8 \to 16 \to -32$ \\
\bottomrule
\end{tabular}

\medskip
\textbf{Check the divergent case.} At $\eta = 1.5$ and $x = 4$:
\[
  x \leftarrow 4 - 1.5(8) = 4 - 12 = -8 ,
\]
then $-8 - 1.5(-16) = -8 + 24 = 16$. Growing by a factor of 2 each step, away
from the minimum. \checkmark

\medskip
\textbf{The stability threshold.} Convergence needs $|1-2\eta| < 1$, so
$0 < \eta < 1$. In general the bound is $\eta < 2/\lambda_{\max}$ where
$\lambda_{\max}$ is the largest Hessian eigenvalue. Here $f'' = 2$, giving
$\eta < 1$. \checkmark

\medskip
\textbf{The conditioning trap.} For $f(x,y) = x^2 + 3y^2$, stability demands
$\eta < 2/6 = 0.333$ from the $y$ direction, while the $x$ direction would happily
take $\eta$ up to $1$. The steepest direction caps the rate for every direction,
so the shallow one is starved.
}

\formaldef{
The learning rate $\eta$ must satisfy $\eta < 2/\lambda_{\max}(H)$ for stability
on a quadratic. Common schedules reduce it over training:
\begin{itemize}
  \item \textbf{step decay}: multiply by a factor at fixed intervals;
  \item \textbf{cosine annealing}: decay smoothly to near zero;
  \item \textbf{warmup}: start small and increase, to avoid early instability.
\end{itemize}
\textbf{Adam} (Kingma \& Ba 2015) instead maintains per-parameter step sizes from
running estimates of the gradient's first and second moments, which adapts to
differing curvature without forming the Hessian (\S\ref{sec:hessian}).
}

\analogybreak{
There is no learning rate that is right throughout training. Early on the surface
is steep and small steps are safe; later it flattens and larger steps would help.
Adaptive methods address this per parameter, and they are heuristics with no
convergence guarantee on non-convex problems --- they work well in practice, which
is a different claim from working provably.
}

\whyhere{
This project uses Adam throughout: $10^{-3}$ for TimeGAN, and $10^{-4}$ for both
generator and critic in QuantGAN and FinGAN. The lower adversarial rate is
standard for WGAN-GP, where a critic that moves too fast destabilises the game of
\S\ref{sec:minimax}. Adam's per-parameter scaling is the practical answer to the
conditioning trap above.
}

%% ─────────────────────────────────────────────────────────────────────────────
\section{Stochastic Gradient Descent and Batching}
\label{sec:sgd}

\situation{
You want the average height of a city's population. You could measure everyone ---
exact, and absurdly slow. Or measure a hundred people, get a good estimate in a
minute, and repeat with a fresh hundred whenever you need another. If you only
need the direction of a trend, the sample is enough.
}

\problem{
The true gradient averages over the whole dataset, and computing it for one update
is wasteful when a small sample gives nearly the same direction. But the sample
introduces noise, and the size of the sample controls a trade-off that this
project got badly wrong once.
}

\workedarith{
\textbf{The estimate and its noise.} With $N$ training windows, the full gradient
is $\frac{1}{N}\sum_{i=1}^{N}\nabla\ell_i$. A batch of $B$ drawn at random gives an
unbiased estimate whose standard error scales as $1/\sqrt{B}$
(\S\ref{sec:lln}).

Quadrupling the batch halves the gradient noise --- and quadruples the cost per
step. Four small steps usually beat one large one.

\medskip
\textbf{What went wrong in this project.} TimeGAN's training loop drew
\emph{one} batch per epoch instead of iterating over all of them.

With $n_{\text{windows}} = 3{,}835$ and $\texttt{batch\_size} = 128$:
\[
  n_{\text{batches}} = \left\lfloor \frac{3835}{128} \right\rfloor = 29 .
\]
So each ``epoch'' performed 1 gradient step where it should have performed 29 ---
a $29\times$ shortfall. Set against QuantGAN, which iterated properly over roughly
59 batches per epoch, TimeGAN was receiving a small fraction of the training its
configuration appeared to specify.

\medskip
\textbf{Why it was invisible.} Both models reported ``epochs'', and the epoch
counts looked comparable. The unit was the same word and meant different things.

\medskip
\textbf{The fix, and then the real fix.} Adding the inner batch loop brought
TimeGAN to 29 steps per epoch. But the deeper problem remained: an epoch is
$\lfloor n/B \rfloor$ steps, so identical epoch counts still mean different step
counts whenever the models or folds differ in size. The project therefore
specifies budgets in \textbf{gradient steps} directly, deriving
\[
  n_{\text{epochs}} = \left\lceil \frac{\text{steps}}{n_{\text{batches}}} \right\rceil
\]
at run time. With a 9{,}000-step budget and 29 batches per epoch that is 311
epochs --- a number nobody would have chosen by hand, and exactly the number
required for parity.
}

\formaldef{
\textbf{Stochastic gradient descent} replaces the full gradient with a mini-batch
estimate:
\[
  \theta_{t+1} = \theta_t - \eta\,\frac{1}{B}\sum_{i \in \mathcal{B}_t}\nabla\ell_i(\theta_t) ,
\]
where $\mathcal{B}_t$ is a randomly drawn batch of size $B$. The estimate is
unbiased: $\E[\hat{g}] = \nabla L$.

An \textbf{epoch} is one pass through the dataset, comprising
$\lfloor N/B \rfloor$ steps. Convergence for SGD requires a decaying learning
rate, $\sum_t \eta_t = \infty$ and $\sum_t \eta_t^2 < \infty$, though constant
rates are used in practice.
}

\analogybreak{
Gradient noise is not purely a cost. It helps escape saddle points and sharp
narrow minima, and models trained with small batches often generalise better than
those trained with very large ones. The noise is part of why the method works, not
only an approximation error to be minimised --- so ``use the largest batch that
fits'' is bad advice.
}

\whyhere{
The epoch-versus-step distinction is the single most consequential lesson in this
project's training code. Reporting a budget in epochs conceals a dependence on
dataset size and batch size; reporting it in gradient steps does not. The parity
assertion in the pipeline now enforces equal \emph{generator updates} across all
three models and refuses to run otherwise --- because a benchmark in which one
model silently trains $29\times$ less than another measures the bug, not the
models.
}

%% ─────────────────────────────────────────────────────────────────────────────
\section{Constraints and Lagrange Multipliers}
\label{sec:lagrange}

\situation{
Fence off the largest possible rectangular field with 100 metres of fencing. You
are not free to make it arbitrarily large --- the fencing is fixed, and the answer
lives on that boundary rather than at an unconstrained optimum.
}

\problem{
Setting the gradient to zero finds unconstrained optima. When the answer is forced
onto a constraint surface, the gradient there is generally \emph{not} zero, and a
different condition is needed.
}

\workedarith{
Maximise area $A = xy$ subject to perimeter $2x + 2y = 100$, i.e.\ $x + y = 50$.

\medskip
\textbf{By substitution.} $y = 50 - x$, so $A(x) = x(50-x) = 50x - x^2$. Then
$A'(x) = 50 - 2x = 0$ gives $x = 25$, $y = 25$, $A = 625$. A square. And
$A'' = -2 < 0$ confirms a maximum. \checkmark

\medskip
\textbf{By Lagrange.} Form
\[
  \mathcal{L}(x,y,\lambda) = xy - \lambda(x + y - 50) .
\]
Setting the partial derivatives to zero:
\[
  \frac{\partial\mathcal{L}}{\partial x} = y - \lambda = 0 , \qquad
  \frac{\partial\mathcal{L}}{\partial y} = x - \lambda = 0 , \qquad
  \frac{\partial\mathcal{L}}{\partial \lambda} = -(x + y - 50) = 0 .
\]
The first two give $x = y = \lambda$; the third gives $2\lambda = 50$, so
$\lambda = 25$ and $x = y = 25$. \checkmark Same answer.

\medskip
\textbf{What $\lambda$ means.} It is the rate at which the optimum improves as the
constraint is relaxed. With 102 metres of fence the optimum is $25.5^2 = 650.25$,
up $25.25$ from $625$ for 2 extra metres --- about $12.6$ per metre, and
$\lambda/2 = 12.5$ per metre of fence since perimeter uses two metres per unit of
$x+y$. \checkmark The multiplier is a shadow price.
}

\formaldef{
To optimise $f(\mathbf{x})$ subject to $g(\mathbf{x}) = 0$, form
\[
  \mathcal{L}(\mathbf{x}, \lambda) = f(\mathbf{x}) - \lambda\,g(\mathbf{x})
\]
and solve $\nabla_{\mathbf{x}}\mathcal{L} = \mathbf{0}$, $g(\mathbf{x}) = 0$.
Geometrically this requires $\nabla f = \lambda\nabla g$: the gradients are
parallel, so no movement along the constraint surface improves $f$.

For inequality constraints $g(\mathbf{x}) \leq 0$ the \textbf{KKT conditions}
apply, adding $\lambda \geq 0$ and complementary slackness $\lambda g = 0$ --- the
multiplier is non-zero only where the constraint binds.
}

\analogybreak{
Exact constraint enforcement is rare in deep learning, because projecting onto a
constraint surface after every step is expensive and often ill-defined. The usual
substitute is a \textbf{penalty}: add a term that grows when the constraint is
violated, and accept approximate satisfaction. That is a soft version of the same
idea with a fixed multiplier chosen by hand rather than solved for.
}

\whyhere{
The WGAN-GP gradient penalty is precisely such a soft constraint. The critic
should satisfy $\norm{\nabla D}_2 = 1$; rather than enforce it exactly, the
objective adds $\lambda\,\E[(\norm{\nabla D}_2 - 1)^2]$ with $\lambda = 10$ fixed.
The constraint is encouraged, never guaranteed --- a distinction
Chapter~\ref{ch:gans} takes seriously.
}

%% ─────────────────────────────────────────────────────────────────────────────
\section{Minimax Problems}
\label{sec:minimax}

\situation{
Rock--paper--scissors. Whatever you settle on, your opponent has a best reply, and
against that reply you would rather have chosen differently. There is no pair of
fixed choices both players are content with, so no amount of thinking converges to
one.
}

\problem{
Everything so far assumed a single objective going downhill. A GAN has two
networks with opposed objectives, and the surface each descends is reshaped by the
other's movement. The convergence story from \S\ref{sec:gradientdescent} does not
transfer, and understanding why is the point of this section.
}

\workedarith{
\textbf{The simplest adversarial objective.} Let
\[
  L(x, y) = xy ,
\]
with one player minimising over $x$ and the other maximising over $y$. The
equilibrium is $(0,0)$: neither can improve unilaterally.

Simultaneous gradient steps with rate $\eta$:
\[
  x \leftarrow x - \eta y , \qquad y \leftarrow y + \eta x .
\]

Start at $(1, 0)$ with $\eta = 0.1$:

\medskip
\begin{tabular}{@{}lrrr@{}}
\toprule
Step & $x$ & $y$ & $\sqrt{x^2+y^2}$ \\
\midrule
0 & $1.000$ & $0.000$ & $1.000$ \\
1 & $1.000$ & $0.100$ & $1.005$ \\
2 & $0.990$ & $0.200$ & $1.010$ \\
3 & $0.970$ & $0.299$ & $1.015$ \\
4 & $0.940$ & $0.396$ & $1.020$ \\
\bottomrule
\end{tabular}

\medskip
\textbf{The distance from equilibrium grows.} It should be shrinking. Compute it
exactly: after one joint step,
\[
  x'^2 + y'^2 = (x - \eta y)^2 + (y + \eta x)^2 = (1 + \eta^2)(x^2 + y^2) .
\]
The radius multiplies by $\sqrt{1+\eta^2} > 1$ at every step, \emph{for any
positive $\eta$}. The players spiral outward for ever.

\medskip
\textbf{This is not a tuning failure.} No learning rate fixes it --- smaller $\eta$
only slows the divergence. The problem is structural: each player descends
correctly with respect to a target the other is moving.

\medskip
\textbf{Contrast with minimisation.} For $f(x) = x^2$ the same simultaneous update
gave a factor of $0.8$ per step (\S\ref{sec:gradientdescent}), monotonically
decreasing. Minimisation has a Lyapunov function --- the objective itself, which
falls every step. A minimax game has none.
}

\formaldef{
A \textbf{minimax problem} seeks
\[
  \min_{\theta}\max_{\phi} V(\theta, \phi) .
\]
A \textbf{Nash equilibrium} is a pair $(\theta^{*},\phi^{*})$ at which neither
player can improve alone:
\[
  V(\theta^{*},\phi^{*}) \leq V(\theta,\phi^{*}) \ \forall\theta ,
  \qquad
  V(\theta^{*},\phi^{*}) \geq V(\theta^{*},\phi) \ \forall\phi .
\]
For convex--concave $V$ an equilibrium exists and averaged iterates converge. For
the non-convex, non-concave $V$ of a GAN, neither existence nor convergence is
guaranteed, and simultaneous gradient steps can cycle or diverge as above.
}

\analogybreak{
The word ``converged'' does not mean for a GAN what it means for a regression. GAN
losses characteristically oscillate without settling, and a stable loss can
indicate collapse rather than success --- a generator producing one output makes
the critic's job easy and both losses flat. Loss curves are close to uninformative
as a stopping criterion, which is why this project judges the models by external
evaluation and explicit degeneracy guards instead.
}

\whyhere{
This is the structural reason GAN training is hard, and the reason the pipeline
carries the safeguards it does: a post-generation check on mean, standard
deviation and ACF(1); a fixed gradient-step budget rather than training to
convergence; and a full evaluation framework external to the training objective.
Chapter~\ref{ch:gans} shows what the failures look like in practice, including a
TimeGAN collapse to standard deviation $0.0002$ against a real $0.0103$.
}
